\section{The Bailout Protocol}\label{sec:bailout}

The Bailout Protocol is the mandatory exception-escalation mechanism of the \bxthree{} Framework. It is not an error-handling routine selected at runtime, nor an optional escalation pathway available to a system that chooses to invoke it. It is an architectural compulsion: any unresolved exception state encountered by any node in a \bxthree{} system \emph{must} propagate upward through the accountability chain until it reaches a named human principal, and no machine actor along that chain may absorb, resolve, or suppress the exception.

This section provides the complete formal specification: the deterministic state machine governing exception propagation (\S~\ref{sec:bailout-state-machine}), the precise conditions governing each transition (\S~\ref{sec:bailout-transitions}), the bail-out trigger condition (\S~\ref{sec:bailout-trigger}), the escalation algorithm (\S~\ref{sec:bailout-escalation-algorithm}), the bypass mechanism enforcing the no-machine-actor property (\S~\ref{sec:bailout-bypass}), and the timeout handling that guarantees liveness under all failure conditions (\S~\ref{sec:bailout-timeouts}).

The protocol operates against the backdrop of the \bxthree{} three-layer model. The \purpose{} layer carries intent, judgment, and final authority; it defines the operating range $R(n)$ for each node $n$ and receives all escalated exceptions. The \bounds{} layer carries bounded reasoning and constrained execution within $R(n)$; it evaluates whether proposed actions fall within the node's authority and initiates bailout when they do not. The \fact{} layer carries deterministic enforcement and forensic auditability; it verifies proposed actions against the authorization ledger, enforces the state machine's transition relation, and maintains the append-only Forensic Ledger recording every protocol event as an immutable architectural record. Each spawned node carries an immutable parent pointer $\alpha(n)$, and every node carries a reference to its human accountability anchor $h(n)$ --- a named human principal or institutional actor --- established at node provisioning and never modified for the node's lifetime.

% ---------------------------------------------------------------
\subsection{State Machine}\label{sec:bailout-state-machine}
% ---------------------------------------------------------------

The Bailout Protocol is formally modeled as a deterministic finite state machine
\begin{equation}
\mathcal{B} \;=\; (Q,\; q_0,\; \Sigma,\; \rightarrow,\; Q_F) \tag{BM}
\end{equation}
where:

\begin{itemize}\itemsep=0pt
\item[$Q$] $\{\;\text{IDLE},\; \text{BOUNDED},\; \text{BAIL\_PENDING},\; \text{ESCALATING},\; \text{HUMAN\_ACKNOWLEDGED},\; \text{RESOLVED}\;\}$ is the finite set of protocol states.
\item[$q_0$] $\text{IDLE}$ is the designated initial state, entered at node startup and after every successful directive resolution.
\item[$\Sigma$] The input alphabet $\{e_{\text{trigger}},\; e_{\text{auth}},\; e_{\text{timeout}},\; e_{\text{ack}},\; e_{\text{resolve}},\; e_{\text{reject}}\}$ comprising trigger events, authorization confirmations, timeout signals, human acknowledgments, binding resolution directives, and rejection directives.
\item[$\rightarrow\; \subseteq Q \times \Sigma \times Q$] The deterministic transition relation defined in Table~\ref{tab:bailout-transitions}.
\item[$Q_F$] $\{\text{RESOLVED}\}$ is the set of terminal accepting states.
\end{itemize}

A node's state $q \in Q$ is stored in the node's Fact Layer worksheet and is observable by its parent node and by the Bailout Monitor. State updates are atomic: no intermediate or partially-executing states are observable externally. The machine is \emph{deterministic}: for every $(q, \sigma) \in Q \times \Sigma$, at most one $q' \in Q$ satisfies $(q, \sigma, q') \in \rightarrow$. Determinism is enforced by the Bailout Monitor's priority function $\pi$, which totally orders simultaneous trigger conditions and selects the highest-priority enabled transition before committing any state change. If no transition is defined for a $(q, \sigma)$ pair, the machine stalls and the Bailout Monitor records a protocol violation in the Forensic Ledger.

\begin{table}[htbp]
\centering
\caption{Bailout Protocol state transition relation $\rightarrow \;\subseteq Q \times \Sigma \times Q$. Rows are current state; columns are input events. Entries denote the next state. Blank cells indicate no transition defined --- the machine stalls and a protocol violation is recorded.}
\label{tab:bailout-transitions}
\begin{tabular}{@{}lcccccc@{}}
\toprule
 & \multicolumn{6}{c}{Input event} \\
\cmidrule(l){2-7}
Current state & $e_{\text{trigger}}$ & $e_{\text{auth}}$ & $e_{\text{timeout}}$ & $e_{\text{ack}}$ & $e_{\text{resolve}}$ & $e_{\text{reject}}$ \\
\midrule
IDLE                  & BOUNDED               & --- & --- & --- & --- & --- \\
BOUNDED               & --- & HUMAN\_ACKNOWLEDGED & BAIL\_PENDING & --- & --- & --- \\
BAIL\_PENDING         & --- & --- & ESCALATING & --- & --- & --- \\
ESCALATING            & --- & HUMAN\_ACKNOWLEDGED & ESCALATING & IDLE & RESOLVED & RESOLVED \\
HUMAN\_ACKNOWLEDGED   & --- & --- & --- & IDLE & RESOLVED & IDLE \\
RESOLVED              & --- & --- & --- & --- & --- & --- \\
\bottomrule
\end{tabular}
\end{table}

\bigskip
\noindent\textbf{State definitions.}

\begin{itemize}
\item[\textbf{IDLE.}] The node is operating within its provisioned operating range $R(n)$. No exception condition is active. The \fact{} layer dispatches normal action requests through the standard execution path. The node transitions to BOUNDED upon receipt of $e_{\text{trigger}}$ satisfying the bail-out trigger condition (TC), defined in \S~\ref{sec:bailout-trigger}.

\item[\textbf{BOUNDED.}] An exception $x$ has been detected by the \bounds{} layer and is being held at the \bounds{}--\fact{} layer boundary. A time-bounded resolution attempt is in progress within the current node's authority. All resolution paths within $R(n)$ remain open. If the Bounds Engine produces a resolution $r$ that the \fact{} layer approves before $T_{\text{bounds}}$ expires, the node returns to IDLE via $e_{\text{auth}}$ (T2). If all resolution paths are exhausted or $T_{\text{bounds}}$ expires, the state advances to BAIL\_PENDING via $e_{\text{timeout}}$ (T3).

\item[\textbf{BAIL\_PENDING.}] The node has determined --- either through an explicit \texttt{bailout\_signal} from the Bounds Engine or through $T_{\text{bounds}}$ expiration without a \fact-layer-approved resolution --- that exception $x$ exceeds its current authority. The Bailout Protocol is formally activated, but propagation to $h(n)$ has not yet been initiated. This is the last architectural point at which a machine actor could theoretically suppress escalation. It cannot: entry to BAIL\_PENDING triggers an immediate atomic transition to ESCALATING via $e_{\text{timeout}}$ (T4), with no dwell time and no machine-actor intervention permitted.

\item[\textbf{ESCALATING.}] The exception $x$ is propagating upward toward $h(n)$ along the parent chain $\alpha^*(\cdot)$. Each intermediate node receives $x$, appends its diagnostic context to the propagation chain (detailed in \S~\ref{sec:bailout-escalation-algorithm}), and forwards the exception via \texttt{SendToParent} without attempting local resolution. The state remains ESCALATING until either (a) $h(n)$ acknowledges receipt ($e_{\text{ack}}$, T5), (b) $h(n)$ issues a binding directive ($e_{\text{resolve}}$ or $e_{\text{reject}}$, T7), or (c) all retry pathways exhaust after $N_{\max}$ retries at $T_{\text{cap}}$, propagating \texttt{PATHWAY\_EXHAUSTED} to $h(n)$ via an alternate functional pathway.

\item[\textbf{HUMAN\_ACKNOWLEDGED.}] $h(n)$ has received $x$, observed the diagnostic context, and acknowledged it. The node awaits a directive. No machine actor may issue, simulate, or substitute a directive in this state. The node remains until $h(n)$ issues \texttt{ACK} ($e_{\text{ack}}$, T8 --- return to IDLE), or a binding directive \texttt{RESOLVE} / \texttt{REJECT} ($e_{\text{resolve}}$ or $e_{\text{reject}}$, T9 --- enter RESOLVED).

\item[\textbf{RESOLVED.}] $h(n)$ has issued a binding directive and the protocol has recorded it in the authorization ledger. The directive is one of: \texttt{ACK} (continue with current parameters, return to IDLE), \texttt{RESOLVE} (execute a specific binding resolution, enter quiescent resolved state), or \texttt{REJECT} (disallow the proposed action, enter quiescent error state). This state is terminal for the current protocol run. On \texttt{ACK}, the node returns to IDLE; on \texttt{REJECT}, the node becomes quiescent, awaiting human re-engagement.
\end{itemize}

\bigskip
\noindent\textbf{State invariants.} The following invariants hold for all reachable states of $\mathcal{B}$:

\begin{align}
\text{INV}_1 &: \forall n \in \text{nodes},\; \text{state}(n) \in Q \setminus \{\text{ESCALATING}\} \iff \neg\text{active\_bailout}(n). \tag{INV-1} \\
\text{INV}_2 &: \text{state}(n) = \text{BAIL\_PENDING} \implies \text{active\_bailout}(n) \land \text{bailout\_reason}(n) \neq \emptyset. \tag{INV-2} \\
\text{INV}_3 &: \text{state}(n) = \text{ESCALATING} \implies \text{active\_bailout}(n) \land \text{propagation\_path}(n) \text{ is defined}. \tag{INV-3} \\
\text{INV}_4 &: \text{state}(n) = \text{RESOLVED} \iff \text{directive}(n) \in \{\text{ACK}, \text{RESOLVE}, \text{REJECT}\}. \tag{INV-4}
\end{align}

INV-1 establishes that non-escalating states are precisely those with no active bailout in flight. INV-2 ensures that BAIL\_PENDING is always accompanied by a non-empty bailout reason. INV-3 ensures that ESCALATING always has a defined propagation path. INV-4 establishes that RESOLVED is terminal and always accompanied by a directive. Together, the invariants enforce that a protocol run cannot enter a state from which no defined transition leads forward.

% ---------------------------------------------------------------
\subsection{Transition Conditions}\label{sec:bailout-transitions}
% ---------------------------------------------------------------

Each transition $T_i$ in the state machine is associated with explicit preconditions that must be satisfied before the transition fires. Transitions are atomic: a transition either commits fully or does not commit at all; there is no partially-executing transition observable in the Forensic Ledger.

\bigskip
\noindent\textbf{Transition table.}

\begin{align}
\text{T1} &: \text{IDLE} \xrightarrow{e_{\text{trigger}}} \text{BOUNDED} && \text{Pre: } \text{TRIGGER}(n, x) = \text{TRUE} \tag{T1} \\
\text{T2} &: \text{BOUNDED} \xrightarrow{e_{\text{auth}}} \text{IDLE} && \text{Pre: } \text{FactLayer approves resolution } r \text{ before } T_{\text{bounds}} \tag{T2} \\
\text{T3} &: \text{BOUNDED} \xrightarrow{e_{\text{timeout}}} \text{BAIL\_PENDING} && \text{Pre: } T_{\text{bounds}} \text{ expired without FactLayer-approved } r \tag{T3} \\
\text{T4} &: \text{BAIL\_PENDING} \xrightarrow{e_{\text{timeout}}} \text{ESCALATING} && \text{Pre: } \text{entry to BAIL\_PENDING (no dwell, no gate)} \tag{T4} \\
\text{T5} &: \text{ESCALATING} \xrightarrow{e_{\text{ack}}} \text{HUMAN\_ACKNOWLEDGED} && \text{Pre: } h(n) \text{ confirmed receipt of exception} \tag{T5} \\
\text{T6} &: \text{ESCALATING} \xrightarrow{e_{\text{timeout}}} \text{ESCALATING} && \text{Pre: } T_{\text{prop}} \text{ expired, } r < N_{\max} \tag{T6} \\
\text{T7} &: \text{ESCALATING} \xrightarrow{e_{\text{resolve}} \cup e_{\text{reject}}} \text{RESOLVED} && \text{Pre: } h(n) \text{ issued binding directive} \tag{T7} \\
\text{T8} &: \text{HUMAN\_ACKNOWLEDGED} \xrightarrow{e_{\text{ack}}} \text{IDLE} && \text{Pre: } h(n) \text{ issued ACK} \tag{T8} \\
\text{T9} &: \text{HUMAN\_ACKNOWLEDGED} \xrightarrow{e_{\text{resolve}} \cup e_{\text{reject}}} \text{RESOLVED} && \text{Pre: } h(n) \text{ issued RESOLVE or REJECT} \tag{T9}
\end{align}

\bigskip
\noindent\textbf{Transition T4: No dwell property.} The most architecturally significant transition is T4. Its preconditions are satisfied immediately upon entry to BAIL\_PENDING; there is no conditional test between T3 firing and T4 firing. This is the \emph{no-dwell property}:
\begin{equation}
\forall n \in \text{nodes},\; \text{state}(n) = \text{BAIL\_PENDING} \implies \text{next\_event}(n) = e_{\text{timeout}} \tag{ND}
\end{equation}
The no-dwell property is enforced by the Fact Layer pre-commit hook: the Bailout Monitor atomically commits T3 and T4 as a single unit, with no opportunity for any machine actor to insert a suppression step between them. This eliminates the last possible intervention point for a machine actor seeking to absorb the exception.

% ---------------------------------------------------------------
\subsection{Bail-Out Trigger Condition}\label{sec:bailout-trigger}
% ---------------------------------------------------------------

The bail-out trigger is the condition that causes a node to enter BOUNDED state (T1). Rather than requiring system designers to enumerate specific error codes in advance --- a practice that provides coverage only for anticipated failure modes and leaves unanticipated ones invisible --- the Bailout Protocol defines the trigger in terms of \emph{observational boundary violation}.

\bigskip
\noindent\textbf{Trigger condition (TC).} A node $n$ triggers bailout when its observable input space contains a condition $x$ such that:
\begin{equation}
x \notin R(n) \;\land\; \neg\text{FactLayer.can\_resolve}(n, x, T_{\text{bounds}}) \tag{TC}
\end{equation}
where $R(n)$ is the provisioned operating range of node $n$, and $\text{FactLayer.can\_resolve}(n, x, T_{\text{bounds}})$ returns \texttt{TRUE} if and only if the \fact{} layer can produce an approved resolution for $x$ within the bounded resolution window $T_{\text{bounds}}$ using only actions within $n$'s current authority. The trigger fires on any condition outside $R(n)$, regardless of whether the condition was anticipated at node provisioning.

\bigskip
\noindent\textbf{Operating range definition.} $R(n)$ is defined at node provisioning by the \purpose{} layer and is stored in the authorization ledger as a declarative specification. It is not a list of permitted actions; it is a constraint region in the node's input space. Formally:
\begin{equation}
R(n) \;=\; \{\; x \in \text{inputSpace}(n) \mid \text{BoundsEngine}(n, x) = \text{WITHIN\_SCOPE} \;\} \tag{OR}
\end{equation}
Conditions within $R(n)$ are handled by the standard execution path. Conditions outside $R(n)$ are candidate exceptions that may require escalation.

\bigskip
\noindent\textbf{Trigger priority function.} When multiple trigger conditions fire simultaneously for node $n$, they are totally ordered by the priority function $\pi(n, x)$, which ensures deterministic transition selection:
\begin{equation}
\pi(n, x) \;=\; w_1 \cdot \text{severity}(x) \;+\; w_2 \cdot \bigl(1 - \text{confidence}(n, x)\bigr) \;+\; w_3 \cdot \text{distance}\bigl(n, h(n)\bigr) \tag{PF}
\end{equation}
where $\text{severity}(x) \in [0, 1]$ is the assigned severity of exception $x$ (higher values indicate greater potential impact), $\text{confidence}(n, x) \in [0, 1]$ is the node's confidence that it can produce a correct resolution for $x$ without human input, $\text{distance}(n, h(n))$ is the number of hops from $n$ to the human anchor in the parent chain, and $w_1 > w_2 > w_3 \geq 0$ are provisioned constants. The highest-priority trigger fires first. Ordering is enforced by the Bailout Monitor before any state transition is committed to the Forensic Ledger.

\bigskip
\noindent\textbf{Secondary trigger: explicit \texttt{bailout\_signal}.} In addition to TC, the Bounds Engine may emit an explicit \texttt{bailout\_signal} at any time, regardless of confidence scores. This handles conditions --- for example, a logical contradiction between two inputs both within $R(n)$ --- that the Bounds Engine recognizes as unresolvable even though each input is individually within the operating range:
\begin{equation}
\texttt{bailout\_signal}(x) \;\implies\; \text{T1 fires for node } n \text{ with exception } x. \tag{TC-secondary}
\end{equation}
The \texttt{bailout\_signal} bypasses TC entirely and immediately initiates the state machine at BOUNDED, where the standard resolution window $T_{\text{bounds}}$ applies before escalation.

\bigskip
\noindent\textbf{The no-discretion property.} The \bounds{} layer does not exercise discretion over the trigger condition. Specifically: (1) The evaluation of $\text{TRIGGER}(n, x)$ is a read-only query against the operating range definition $R(n)$, not a judgment call made by a machine actor. (2) The result of the query directly drives the state machine transition without an intervening decision step. (3) The confidence function and severity weights are provisioned at node initialization and are immutable at runtime. These constraints ensure that the trigger condition is architecturally enforced rather than operationally optional.

% ---------------------------------------------------------------
\subsection{Escalation Algorithm}\label{sec:bailout-escalation-algorithm}
% ---------------------------------------------------------------

When the state machine enters ESCALATING (T4), the protocol executes the \texttt{Escalate} algorithm to propagate the exception to $h(n)$ along the parent chain. The algorithm is executed by the Bailout Monitor on the originating node; intermediate nodes execute only \texttt{ForwardException}, which appends diagnostic context and relays the exception upward without attempting local resolution.

\begin{algorithm}[htbp]
\caption{Bailout Protocol escalation algorithm.}
\label{alg:escalate}
\begin{algorithmic}[1]
\Procedure{Escalate}{$n$, trigger}
  \State $\text{ctx} \leftarrow$ \Call{BuildDiagnosticContext}{$n$, trigger}
  \State $p \leftarrow \alpha(n)$ \Comment{immediate parent node}
  \State $\text{path} \leftarrow$ \Call{SelectPathway}{PHR}
  \State $r \leftarrow 0$ \Comment{retry counter}
  \While{$\text{true}$}
    \State $\text{sent} \leftarrow$ \Call{SendToParent}{$p$, ctx, path}
    \If{$\text{sent.acknowledged}$}
      \State \Return \texttt{ESCALATING\_CONFIRMED}
    \EndIf
    \State $r \leftarrow r + 1$
    \If{$r \geq N_{\max}$}
      \State \Call{RecordPathwayExhausted}{$n$, ctx}
      \State \Call{PropagateToHumanRoot}{$n$, ctx, trigger}
      \State \Return \texttt{PATHWAY\_EXHAUSTED}
    \EndIf
    \State $T_{\text{prop}} \leftarrow \min(2 \cdot T_{\text{prop}},\; T_{\text{cap}})$
    \State \Call{Wait}{$T_{\text{prop}}$}
    \State $\text{path} \leftarrow$ \Call{SelectPathway}{PHR}
  \EndWhile
\EndProcedure
\\[4pt]
\Procedure{ForwardException}{$n$, ctx}
  \State $\text{ctx}' \leftarrow \text{ctx}$
  \State $\text{ctx}'.\text{diagnosticChain} \leftarrow \text{ctx}'.\text{diagnosticChain} \;\;[n.\text{nodeId}, n.\text{timestamp}, n.\text{local\_error\_summary}]$
  \State $p \leftarrow \alpha(n)$
  \State \Call{SendToParent}{$p$, $\text{ctx}'$, \Call{SelectPathway}{PHR}}
\EndProcedure
\\[4pt]
\Procedure{SelectPathway}{PHR}
  \State \Return $\text{path} \in \text{PHR.pathways}$ \text{ with minimum estimated latency among } \text{path.status} = \texttt{FUNCTIONAL}
\EndProcedure
\end{algorithmic}
\end{algorithm}

\bigskip
\noindent\textbf{BuildDiagnosticContext.} The diagnostic context $\text{ctx}$ is the data structure that travels unchanged through each level of the parent chain until it reaches $h(n)$. It comprises:
\begin{align}
\text{ctx} \;=\; \bigl(&\text{exceptionId},\; \text{originatingNode},\; \text{bailoutReason},\; \text{triggerCondition}, \bigr. \tag{DC} \\
  &\bigl.\text{operatingRange},\; \text{attemptedResolutions},\; \text{diagnosticChain},\; \text{timestamp},\; \text{severity}\bigr)
\end{align}
where $\text{diagnosticChain}$ is appended to at each hop by \texttt{ForwardException}, creating a temporally ordered record of every node that handled the exception and the local error summary produced at each level. This chain is the primary forensic artifact used by $h(n)$ to understand the full context of the exception.

\bigskip
\noindent\textbf{PropagateToHumanRoot.} When all parent-pathway retries are exhausted ($r \geq N_{\max}$), the algorithm invokes \texttt{PropagateToHumanRoot}, which propagates the exception directly to $h(n)$ via an alternate functional pathway tracked by the Pathway Health Registry. This ensures that pathway exhaustion at the immediate parent level does not block escalation to the human anchor.

% ---------------------------------------------------------------
\subsection{Bypass Mechanism}\label{sec:bailout-bypass}
% ---------------------------------------------------------------

The parent-chain bypass is the property that the Bailout Protocol's escalation pathway contains no machine actors as terminal or intermediate absorption points between the originating node and $h(n)$. Formally:
\begin{equation}
\forall n \in \text{nodes},\; \forall m \in \text{machineActors},\; m \in \alpha^+(n) \implies m \notin \text{termini}(n) \tag{BP}
\end{equation}
where $\alpha^+(n)$ is the transitive closure of the parent relation (the chain from $n$ to the human root), and $\text{termini}(n)$ is the set of permissible termination points for node $n$'s exception propagation. By definition, $\text{termini}(n) = \{h(n)\}$.

\bigskip
\noindent\textbf{Why bypass is architecturally necessary.} The bypass is not an optimization --- it is a structural necessity. Any machine actor $m$ that could absorb an exception in ESCALATING state would become a terminal node for that exception, and the accountability chain would terminate at $m$ rather than at a human. If $m$ then fails, or if $m$'s resolution is incorrect or insufficient, no human principal is present in the chain to detect or correct the failure. The system compounds failure instead of surfacing it.

This is the precise failure mode the Bailout Protocol is designed to prevent. The functional separation between \purpose{}, \bounds{}, and \fact{} layers creates the structural condition that makes bypass necessary: the \bounds{} layer detects exception conditions but carries no judgment authority; the \fact{} layer enforces transition relations and records events but carries no intent authority. When an exception exceeds the \bounds{} layer's provisioned scope, no machine actor exists that has the authority to adjudicate it. The decision must return to the \purpose{} layer, which carries the intent and final accountability. The Bailout Protocol is the mechanism by which this structural necessity is enforced at runtime.

\bigskip
\noindent\textbf{Enforcement mechanisms.} The bypass is enforced by two cooperating mechanisms:

\begin{enumerate}
\item[\textbf{Fact Layer gate.}] The \fact{} layer is the only mechanism capable of authorizing state transitions in the \bxthree{} Framework. The Bailout Monitor's pre-commit hook rejects any ledger entry for a resolution where the underlying trigger satisfies TC and $T_{\text{bounds}}$ has expired. This prevents the Bounds Engine from issuing a retroactive resolution after timeout to suppress warranted escalation. The \fact{} layer gate operates identically on all machine actors: there is no privileged machine actor with the ability to bypass this check.

\item[\textbf{Immutable parent-pointer and human-root architecture.}] The parent pointer $\alpha(n)$ and the human-root identifier $h(n)$ are stored in \fact-layer-managed memory and are read-only at the machine-actor level. A machine actor cannot redirect the parent chain, cannot insert itself as an intermediate resolver, and cannot modify the $humanRootId$ field that terminates the chain. Any attempted modification requires \purpose{} layer authorization with a corresponding ledger entry --- which cannot be completed without human approval if the modification would affect the bailout terminus.
\end{enumerate}

Together, these mechanisms ensure that BP holds for all $n \in \text{nodes}$ under all execution conditions, including partial node failure, network partition, and deliberate attempt by a compromised machine actor to redirect the bailout pathway.

Conventional escalation hierarchies --- including tiered oversight models that permit absorption at intermediate tiers --- make human oversight a high-tier option rather than a compulsory terminus. The Bailout Protocol forbids this absorption by architectural constraint: $h(n)$ is not one option among many --- it is the \emph{only} permissible terminus for every unresolved exception.

% ---------------------------------------------------------------
\subsection{Timeout Handling}\label{sec:bailout-timeouts}
% ---------------------------------------------------------------

Timeouts in the Bailout Protocol serve a liveness function: they prevent the protocol from permanently stalling when a parent node is unreachable, a communication pathway is unresponsive, or the human anchor is temporarily unavailable. Each timeout has a defined scope, a defined scope-exit action, and an escalation path that prevents deadlock.

\bigskip
\noindent\textbf{Timeout definitions.}

\begin{enumerate}
\item[\textbf{$T_{\text{bounds}}$ --- Bounds Engine resolution timeout.}] The maximum time a node may remain in BOUNDED state attempting autonomous resolution. Default: 60 seconds, provisioned per node class according to the expected latency of the node's resolution procedure. If $T_{\text{bounds}}$ expires without a \fact-layer-approved resolution, the node transitions to BAIL\_PENDING $\rightarrow$ ESCALATING via T3 and T4. The \fact{} layer cannot extend $T_{\text{bounds}}$ unilaterally; any extension requires human authorization via a separate ledger entry recorded before $T_{\text{bounds}}$ expires.

\item[\textbf{$T_{\text{prop}}$ --- Per-hop propagation timeout.}] The maximum time node $n$ waits for an acknowledgment from its parent after sending an exception via \texttt{SendToParent}. Default initial value: 5 seconds. If $T_{\text{prop}}$ expires without ACK, the algorithm retries with exponential backoff: $T_{\text{prop}} \leftarrow \min(2 \cdot T_{\text{prop}},\; T_{\text{cap}})$. This backoff continues until either an ACK is received or $N_{\max}$ retries have been exhausted at $T_{\text{cap}}$.

\item[\textbf{$T_{\text{cap}}$ --- Maximum propagation timeout value.}] The ceiling value for the exponential backoff of $T_{\text{prop}}$. Default: 60 seconds. After $N_{\max}$ consecutive retries at $T_{\text{cap}}$ without ACK, the algorithm fires a \texttt{parent\_unreachable} timeout event, which propagates to $h(n)$ via the next available functional pathway tracked by the Pathway Health Registry.

\item[\textbf{$T_{\text{human}}$ --- Human acknowledgment timeout.}] The maximum time $h(n)$ may remain in HUMAN\_ACKNOWLEDGED state without issuing a directive. Default: 300 seconds (5 minutes), configurable per deployment based on the expected human response time for the node's operational domain. If $T_{\text{human}}$ expires, the protocol issues a reminder notification to $h(n)$ and transitions the originating node to RESOLVED with an \texttt{unresolved} tag. The node becomes quiescent: no autonomous retry occurs, no further actions are attempted, and the \texttt{unresolved} tag in the Forensic Ledger signals that human re-engagement is required to clear the state. The protocol does not re-escalate automatically; automatic re-escalation would create a potential denial-of-service condition against $h(n)$.
\end{enumerate}

\begin{table}[htbp]
\centering
\caption{Timeout parameters, their start conditions, and their functional scopes.}
\label{tab:bailout-timeouts}
\begin{tabular}{@{}lcl@{}}
\toprule
\textbf{Timeout} & \textbf{Started at} & \textbf{Scope} \\
\midrule
$T_{\text{bounds}}$ & T1 fires (TC satisfied) & Bounds Engine autonomous resolution window \\
$T_{\text{prop}}$ & T4 fires (transmit to $h(n)$) & Per-hop parent acknowledgment \\
$T_{\text{cap}}$ & T6 fires (retry) & Ceiling for $T_{\text{prop}}$ backoff \\
$T_{\text{human}}$ & ESCALATING + $h(n)$ contact confirmed & Human directive issuance window \\
\bottomrule
\end{tabular}
\end{table}

\bigskip
\noindent\textbf{Pathway Health Registry (PHR).} The PHR is a runtime data structure maintained by the Bailout Monitor that tracks the availability, latency, and error rate of each provisioned communication pathway to $h(n)$. Before each \texttt{SendToParent} call, \texttt{SelectPathway} queries the PHR and selects the pathway with the lowest estimated latency among currently functional pathways. Pathways that fail $K$ consecutive health checks (default $K = 3$) are marked \texttt{DEGRADED} and excluded from pathway selection until a health check succeeds, preventing the algorithm from repeatedly selecting a known-degraded pathway.

The PHR is updated by a background health-ping thread that sends heartbeat messages along each provisioned pathway at intervals of $T_{\text{health}}$ (default: 30 seconds). The health-ping thread operates with the same scheduling priority as the Bailout Monitor and is not preemptible by agent tasks, ensuring that pathway status remains current even under heavy agent workload.

\bigskip
\noindent\textbf{Liveness guarantee.} Together, $T_{\text{bounds}}$, $T_{\text{prop}}$, $T_{\text{cap}}$, and $T_{\text{human}}$ define a finite upper bound on wall-clock time from trigger firing (T1) to human acknowledgment under any combination of parent unavailability and pathway degradation:
\begin{equation}
T_{\max} \;=\; T_{\text{bounds}} \;+\; N_{\max} \cdot T_{\text{cap}} \;+\; T_{\text{human}} \tag{T-max}
\end{equation}

A deployment that cannot guarantee human acknowledgment within the required $T_{\max}$ under all failure scenarios must provision additional or higher-bandwidth pathways, or must reduce $T_{\text{prop}}$ and $T_{\text{cap}}$ to bring the worst-case bound within the required limit. This analysis is a required part of deployment certification for any \bxthree{} system operating in a regulated or safety-critical environment.

\bigskip
\noindent\textbf{Timeout interaction with the bypass mechanism.} The bypass mechanism ensures that timeouts cannot be weaponized by a machine actor to delay or suppress escalation. Specifically: (a) $T_{\text{bounds}}$ is enforced by the \fact{} layer pre-commit hook, not by the \bounds{} layer; a machine actor cannot artificially extend $T_{\text{bounds}}$ to prevent the transition to ESCALATING. (b) $T_{\text{prop}}$ backoff is calculated by the Bailout Monitor, not by any machine actor; the backoff curve is defined in the protocol specification and cannot be altered at runtime. (c) If all pathways time out at $T_{\text{cap}}$, the protocol resolves with \texttt{PATHWAY\_EXHAUSTED} --- this is a terminal RESOLVED state, not a suppression of escalation. The \texttt{PATHWAY\_EXHAUSTED} tag in the Forensic Ledger ensures that the pathway failure is recorded as an incident requiring human remediation, not as a resolved exception.
